% Part: lambda-calculus
% Chapter: introduction
% Section: composition

\documentclass[../../../include/open-logic-section]{subfiles}

\begin{document}

\olfileid{lam}{int}{com}
\olsection{Fungsi kang \usetoken{S}{lambda definable} Katutup tumrap Komposisi}

\begin{lem}
Fungsi-fungsi kang !!{lambda definable}
katutup tumrap komposisi.
\end{lem}

\begin{proof}
Upamana $f$ ditetepake kanthi komposisi
saka $h$, $g_0,$ \dots,~$g_{k-1}$.
Upamana $h$, $g_0$, \dots,~$g_{k-1}$
saben-saben !!{lambda defined} dening
$H$, $G_0$, \dots,~$G_{k-1}$.
Kita kudu nemokake suku~$F$ kang
!!{lambda define}s~$f$. Nanging kita
bisa langsung netepake~$F$ kanthi
\[
F(x_0, \dots, x_{l-1}) = H(G_0(x_0, \dots, x_{l-1}),
\dots, G_{k-1}(x_0, \dots, x_{l-1})).
\]
Tegese, basa kalkulus lambda pancen trep
kanggo makili komposisi.
\end{proof}

\end{document}
